\section{Placa de Medição}
\label{sec:siba}
\label{sec:placa_medicao}
A placa de medição é o nó central do sistema de aquisição: mede as
tensões e correntes de fase do motor e, a partir delas, estima rotação e torque
com o observador; aciona o controlador eletrônico de velocidade (ESC) e os
relés de potência; concentra as verificações de segurança; e coordena a rede de
dados com a interface. Seu desenvolvimento foi dividido em infraestrutura de
\textit{hardware} e lógica embarcada.

\subsection{Microcontrolador e Comunicação}
\label{sec:padronizacao_placas}
Esta subseção estabelece as diretrizes de projeto e as restrições metodológicas adotadas para o desenvolvimento da placa de medição. Visando otimizar os custos de manufatura e garantir a modularidade do sistema, por conta disso irá se usar placas de fenolite manufaturadas a mão, bem como módulos, placas de desenvolvimento e conversores comerciais.

Adotou-se o microcontrolador CH32V307VCT6 (WCH) como unidade de processamento
da placa de medição. A escolha fundamenta-se em quatro critérios principais.
Primeiro, o conversor analógico-digital com taxa de até 1 Msps
atende à amostragem rápida de correntes e tensões exigida pela aquisição e pelo
controle. Segundo, a ampla memória disponível acomoda o \textit{firmware} sobre
o sistema operacional de tempo real Zephyr, as estruturas de dados de aquisição
e o processamento em ponto flutuante do controle vetorial e do observador, este
último amparado pela unidade de ponto flutuante (FPU) do núcleo RISC-V de
\SI{144}{\mega\hertz}. Terceiro, o dispositivo integra interface Ethernet
(MAC e PHY), detalhada adiante. Por fim, o baixo custo — cerca de
R\$30 por placa, ou aproximadamente R\$90 considerando
tributos de importação — viabiliza a produção de múltiplas unidades dentro das
restrições orçamentárias do projeto.

A interface Ethernet merece destaque por ter motivado a escolha. Durante o ano
de 2025, membros da equipe Albatroz Aerodesign identificaram a necessidade de um
método de comunicação mais simples e padronizado entre as placas de aquisição e
os computadores de ensaio. A conectividade por rede via cabo atende a essa
demanda: dispensa conversores de sinal intermediários, como adaptadores
UART-TTL (\textit{Universal Asynchronous Receiver/Transmitter} e \textit{Transistor-Transistor Logic}, e equipamentos dedicados de maior porte, uma vez que a maior parte dos
computadores e \textit{notebooks} dispõe nativamente de conexão de rede cabeada, além de facilidade adicionar futuras demandas não esperadas neste trabalho.

\subsection{Sistema Operacional Vcore}
\label{sec:vcore}

O \textit{firmware} da placa de medição é estruturado em camadas lógicas,
apresentadas na \autoref{fig:Metodologia_Diagrama_placas}. A camada de
aplicação, que nos sistemas embarcados costuma ficar no próprio
microcontrolador, foi deslocada para o computador: as rotinas de ensaio, os
procedimentos de calibração e de estimação, os gráficos e o registro em CSV
são executados pela aplicação de supervisão (\autoref{sec:interface}). A
placa guarda apenas o que precisa responder em tempo real (aquisição,
observador, acionamento e segurança) e recebe da aplicação comandos
simples pela rede. Com isso a memória RAM do microcontrolador, já ocupada em
grande parte pelo Zephyr e pela pilha de rede, fica livre para os
\textit{buffers} de aquisição e para o observador, e mudar um ensaio não exige
regravar o \textit{firmware}.

\begin{figure}[H]
    \centering
    \caption{Camadas lógicas do sistema: aplicação no computador e
    \textit{firmware} da placa de medição}
    \label{fig:Metodologia_Diagrama_placas}
    \begin{tikzpicture}[
        font=\small,
        layer/.style={draw, rounded corners, minimum width=11.6cm, minimum height=1.75cm},
        mod/.style={draw, rounded corners, align=center, minimum height=0.8cm,
            minimum width=2.0cm, fill=white, font=\footnotesize},
        title/.style={font=\bfseries\small},
        arrow/.style={-{Latex}, thick}
    ]
        % ---------------- computador ----------------
        \node[layer, fill=blue!6, dashed, minimum height=1.9cm] (app) at (0,5.0) {};
        \node[title] at (app.north) [yshift=-0.3cm] {Application: aplicação de supervisão (computador)};
        \node[mod] at (-4.45,4.75) {Rotinas de\\ensaio};
        \node[mod] at (-2.22,4.75) {Calibração e\\estimação};
        \node[mod] at (0,4.75) {Ajuste do\\observador};
        \node[mod] at (2.22,4.75) {Gráficos e\\registro CSV};
        \node[mod] at (4.45,4.75) {Perfis de motor\\e hélice};

        \draw[{Latex}-{Latex}, very thick, blue!60!black] (app.south) -- ++(0,-0.75)
            node[midway, right, font=\footnotesize] {Ethernet / UDP: comandos e telemetria};

        % ---------------- microcontrolador ----------------
        \node[draw, thick, rounded corners, minimum width=12.4cm, minimum height=5.3cm,
            fill=black!2] (mcu) at (0,0.45) {};
        \node[font=\footnotesize\itshape, anchor=north west] at (mcu.north west) {microcontrolador CH32V307};

        \node[layer, fill=black!10] (sys) at (0,1.85) {};
        \node[title] at (sys.north) [yshift=-0.3cm] {System: gerenciadores de domínio};
        \node[mod] at (-4.45,1.6) {\textit{Battery}\\\textit{Manager}};
        \node[mod] at (-2.22,1.6) {\textit{Propeller}\\\textit{Manager}};
        \node[mod] at (0,1.6) {\textit{Observer}\\\textit{Manager}};
        \node[mod] at (2.22,1.6) {\textit{Security}\\\textit{Manager}};
        \node[mod] at (4.45,1.6) {\textit{Log}\\\textit{Manager}};

        \node[layer, fill=black!15] (os) at (0,-0.55) {};
        \node[title] at (os.north) [yshift=-0.3cm] {Vcore - OS};
        \node[mod, fill=yellow!15, thick] (ossi) at (-4.45,-0.8) {OSSI\\[-1pt]{\scriptsize fila única}};
        \node[mod] at (-2.22,-0.8) {NCI\\[-1pt]{\scriptsize rede UDP}};
        \node[mod] at (0,-0.8) {HAL\\[-1pt]{\scriptsize periféricos}};
        \node[mod] at (2.22,-0.8) {CMD\_Core\\[-1pt]{\scriptsize comandos}};
        \node[mod] at (4.45,-0.8) {NVM\\[-1pt]{\scriptsize EEPROM}};

        \draw[arrow] (sys.south) -- (os.north) node[midway, right, font=\footnotesize] {despacho via OSSI};
        \node[font=\footnotesize, gray] at (0,-1.6) {Zephyr RTOS e periféricos (ADC + DMA, Ethernet, I\textsuperscript{2}C, PWM)};
    \end{tikzpicture}

    Fonte: Autor (2026)
\end{figure}

A camada base (\textit{Vcore - OS}) encapsula as rotinas nativas do
\textit{Zephyr} e os periféricos do microcontrolador, apresentando uma
arquitetura reaproveitável por outras placas da bancada e promovendo o
reaproveitamento de código entre elas. Ela é composta por cinco módulos, detalhados na
\autoref{fig:Metodologia_Vcore}:

\begin{itemize}
    \item \textbf{OSSI (\textit{Operating System Standard Interface}):} fila de trabalho única do sistema, que recebe os dados publicados pelo HAL e pelo CMD\_Core e despacha o processamento para a camada \textit{System} sempre na mesma tarefa (\autoref{fig:arquitetura_fila_unica}).
    \item \textbf{NCI (\textit{Network Core Interface}):} gerenciador do \textit{socket} UDP, responsável pela recepção e pela transmissão de pacotes com latência minimizada.
    \item \textbf{HAL (\textit{Hardware Abstraction Layer}):} interface direta com os periféricos do microcontrolador (como o ADC operando via DMA para leituras analógicas em \textit{background}), isolando os registradores físicos das camadas superiores.
    \item \textbf{CMD\_Core:} interpretador central que recebe os pacotes binários brutos entregues pelo NCI, valida as instruções e as encaminha ao OSSI para despacho.
    \item \textbf{NVM (\textit{Non-Volatile Memory}):} guarda numa memória EEPROM externa, acessada por barramento I\textsuperscript{2}C, os coeficientes de calibração de cada sensor e os parâmetros do observador. Na partida o módulo restaura o que está gravado; quando o operador confirma uma calibração pela aplicação de supervisão, o CMD\_Core manda gravá-la, de modo que ela sobrevive à desenergização da placa.
\end{itemize}

\begin{figure}[H]
    \centering
    \caption{Módulos da camada \textit{Vcore - OS}}
    \label{fig:Metodologia_Vcore}
    \begin{tikzpicture}[
        font=\small,
        box/.style={draw, rounded corners, align=center, minimum height=0.9cm, minimum width=2.4cm},
        arrow/.style={<->, thick}
    ]
        \node[box] (ossi) {OSSI};
        \node[box, above=1.1cm of ossi] (hal) {HAL};
        \node[box, left=1.6cm of ossi] (nci) {NCI};
        \node[box, right=1.6cm of ossi] (cmd) {CMD\_Core};
        \node[box, below=1.1cm of ossi] (nvm) {NVM\\(EEPROM)};

        \draw[arrow] (hal) -- (ossi);
        \draw[arrow] (nci) -- (ossi);
        \draw[arrow] (cmd) -- (ossi);
        \draw[arrow] (nvm) -- (ossi);
    \end{tikzpicture}

    Fonte: Autor (2026)
\end{figure}

Já a camada \textit{System} concentra os gerenciadores de domínio (bateria,
hélice, segurança, telemetria e o observador SMO), sem
nenhuma \textit{thread} própria: eles são chamados diretamente pela função
de despacho do OSSI, conforme detalhado adiante para a placa de medição
(\autoref{sec:siba}). A camada
\textit{Application} fica no computador, na aplicação de supervisão.

Para evitar que vários eventos alterem as mesmas variáveis ao mesmo tempo, o
processamento da placa é concentrado em uma única fila de trabalho, a do OSSI
(\autoref{fig:arquitetura_fila_unica}). Os eventos externos, como uma amostra
processada, um comando recebido ou o instante de enviar a telemetria, apenas
colocam um item nessa fila, e os gerenciadores de domínio são executados um de
cada vez, na mesma tarefa. Assim, dois gerenciadores nunca disputam o mesmo
estado, e não são necessárias travas entre eles. A única exceção é a
aquisição das fases, que tem fila própria de maior prioridade
(\autoref{sec:same_firmware}).

\begin{figure}[H]
    \centering
    \caption{Arquitetura de fila de trabalho única da placa de medição}
    \label{fig:arquitetura_fila_unica}
    \resizebox{\textwidth}{!}{%
    \begin{tikzpicture}[
        font=\small, >={Latex},
        box/.style={draw, rounded corners, align=center, minimum height=1.1cm, text width=2.8cm, fill=white},
        src/.style={box, fill=black!6},
        ossi/.style={box, fill=yellow!20, thick, text width=2.6cm, minimum height=4.4cm},
        arrow/.style={->, thick},
        rot/.style={font=\scriptsize, fill=white, inner sep=1pt}
    ]
        % fontes de evento e primeiro tratamento
        \node[src] (adc) at (0,1.8)   {Conversor\\(ADC + DMA)};
        \node[box] (acq) at (4.4,1.8) {Fila de aquisição\\{\footnotesize Clarke, SMO e PLL}};
        \node[src] (tmr) at (0,0)     {Temporizador\\(\SI{10}{\hertz})};
        \node[box] (rx)  at (4.4,-1.8) {Recepção (NCI)\\{\footnotesize interpretação (CMD\_Core)}};
        % OSSI
        \node[ossi] (wq) at (8.6,0) {\textbf{OSSI}\\[2pt]fila de\\trabalho única};
        \node[src] (net) at (8.6,-3.9) {Rede\\(UDP)};
        % gerenciadores e envio
        \node[box] (mgr) at (12.8,1.2) {Gerenciadores\\de domínio};
        \node[box] (tx)  at (12.8,-1.2) {\textit{Log Manager}\\{\footnotesize envio UDP (NCI)}};
        \node[src] (pc)  at (17.0,-1.2) {Aplicação\\de supervisão};

        \draw[arrow] (adc) -- (acq);
        \draw[arrow] (acq.east) -- node[rot, above] {resultados} (acq.east -| wq.west);
        \draw[arrow] (tmr.east) -- node[rot, above, pos=0.25] {telemetria} (tmr.east -| wq.west);
        \draw[arrow] (rx.east) -- node[rot, above] {comando} (rx.east -| wq.west);
        \draw[arrow] (net.west) -| (rx.south);
        \draw[arrow] (mgr.west -| wq.east) -- (mgr.west);
        \draw[arrow] (tx.west -| wq.east) -- (tx.west);
        \draw[arrow] (tx) -- (pc);
        \draw[arrow, dashed] (pc.south) |- node[rot, pos=0.75, above] {comandos} (net.east);
    \end{tikzpicture}%
    }

    Fonte: Autor (2026)
\end{figure}

\subsection{\textit{Hardware}}

A placa usa o microcontrolador CH32V307VCT6
(\autoref{sec:padronizacao_placas}) na forma de uma placa de desenvolvimento
comercial, soldada sobre uma placa perfurada que recebe os conectores dos
sensores, dos relés e da alimentação. A \autoref{fig:pinagem_siba} mostra a
pinagem adotada, desenhada na posição real dos dois conectores da placa de
desenvolvimento, que é a referência usada na montagem. A Porta E,
que não tem função analógica, foi reservada para seis células de carga: os
pinos ficam disponíveis na placa e o \textit{firmware} já tem suporte à
leitura, mas a montagem das células e a medição da tração da hélice não fazem
parte do escopo deste trabalho. Os
pinos analógicos das Portas A e C ficam para os dois conversores, em modo
dual simultâneo, e recebem as três correntes e as três tensões de fase vindas
da placa de condicionamento (\autoref{sec:same}).

\begin{figure}[H]
    \centering
    \caption{Pinagem da placa de medição na placa de desenvolvimento do CH32V307VCT6}
    \label{fig:pinagem_siba}
    \resizebox{\textwidth}{!}{% Gerado a partir de "Imgs/Montagem croqui.xlsx" (aba micro) + pinagem do firmware.
\begin{tikzpicture}[font=\fontsize{7}{8}\selectfont, every node/.append style={font=\fontsize{7}{8}\selectfont},
  pin/.style={draw, minimum width=0.56cm, minimum height=0.36cm, inner sep=0pt, font=\fontsize{4.6}{5}\selectfont\sffamily},
  p-pwr/.style={pin, fill=gray!35}, p-free/.style={pin, fill=white, text=gray},
  p-adc/.style={pin, fill=orange!45}, p-nm/.style={pin, fill=orange!15, dashed, text=gray},
  p-dig/.style={pin, fill=cyan!30}, p-com/.style={pin, fill=green!30}, p-seg/.style={pin, fill=red!35}, p-res/.style={pin, fill=white, dashed},
  d/.style={inner sep=1pt}, d-nm/.style={d, text=gray, font=\fontsize{7}{8}\selectfont\itshape}, d-res/.style={d, text=gray}]
  \draw[rounded corners=3pt, very thick, fill=black!4] (-0.38,0.5) rectangle (9.0,-10.16);
  \node[font=\fontsize{9}{10}\selectfont\bfseries] at (4.31,0.85) {SIBA};
  \node[draw, fill=white, minimum width=2.2cm, minimum height=0.7cm, font=\fontsize{5}{6}\selectfont] at (4.31,-9.5) {conector Ethernet};
  \node[draw, fill=white, minimum width=1.7cm, minimum height=1.1cm, font=\fontsize{5}{6}\selectfont, align=center] at (4.6,-2.4) {CH32V307VCT6};
  \node[p-pwr] at (0.00,0.00) {GND};
  \node[p-pwr] at (0.62,0.00) {GND};
  \node[p-pwr] at (0.00,-0.42) {GND};
  \node[p-pwr] at (0.62,-0.42) {GND};
  \node[p-dig] at (0.00,-0.84) {PE0};
  \node[d, anchor=east] at (-0.50,-0.84) {HX711 SCK célula 1};
  \node[p-dig] at (0.62,-0.84) {PE1};
  \node[d, anchor=west] at (0.96,-0.84) {HX711 SCK célula 2};
  \node[p-dig] at (0.00,-1.26) {PE2};
  \node[d, anchor=east] at (-0.50,-1.26) {HX711 SCK célula 3};
  \node[p-dig] at (0.62,-1.26) {PE3};
  \node[d, anchor=west] at (0.96,-1.26) {HX711 SCK célula 4};
  \node[p-dig] at (0.00,-1.68) {PE4};
  \node[d, anchor=east] at (-0.50,-1.68) {HX711 SCK célula 5};
  \node[p-dig] at (0.62,-1.68) {PE5};
  \node[d, anchor=west] at (0.96,-1.68) {HX711 SCK célula 6};
  \node[p-dig] at (0.00,-2.10) {PE6};
  \node[d, anchor=east] at (-0.50,-2.10) {HX711 DOUT célula 1};
  \node[p-free] at (0.62,-2.10) {PC13};
  \node[p-free] at (0.00,-2.52) {PC14};
  \node[p-free] at (0.62,-2.52) {PC15};
  \node[p-free] at (0.00,-2.94) {PC0};
  \node[p-free] at (0.62,-2.94) {PC1};
  \node[p-nm] at (0.00,-3.36) {PC2};
  \node[d-nm, anchor=east] at (-0.50,-3.36) {Reserva (ADC2)};
  \node[p-nm] at (0.62,-3.36) {PC3};
  \node[d-nm, anchor=west] at (0.96,-3.36) {Tensão bateria (ADC2)};
  \node[p-nm] at (0.00,-3.78) {PA0};
  \node[d-nm, anchor=east] at (-0.50,-3.78) {Pitot (ADC1)};
  \node[p-nm] at (0.62,-3.78) {PA1};
  \node[d-nm, anchor=west] at (0.96,-3.78) {Termistor bateria (ADC1)};
  \node[p-res] at (0.00,-4.20) {PA2};
  \node[d-res, anchor=east] at (-0.50,-4.20) {Reserva \textit{flash}};
  \node[p-nm] at (0.62,-4.20) {PA3};
  \node[d-nm, anchor=west] at (0.96,-4.20) {Termistor motor (ADC1)};
  \node[p-free] at (0.00,-4.62) {PA4};
  \node[p-res] at (0.62,-4.62) {PA5};
  \node[d-res, anchor=west] at (0.96,-4.62) {Reserva \textit{flash}};
  \node[p-res] at (0.00,-5.04) {PA6};
  \node[d-res, anchor=east] at (-0.50,-5.04) {Reserva \textit{flash}};
  \node[p-res] at (0.62,-5.04) {PA7};
  \node[d-res, anchor=west] at (0.96,-5.04) {Reserva \textit{flash}};
  \node[p-nm] at (0.00,-5.46) {PC4};
  \node[d-nm, anchor=east] at (-0.50,-5.46) {Corrente bateria (ADC2)};
  \node[p-free] at (0.62,-5.46) {PC5};
  \node[p-dig] at (0.00,-5.88) {PB0};
  \node[d, anchor=east] at (-0.50,-5.88) {RPM (TIM3 CH3)};
  \node[p-free] at (0.62,-5.88) {PB1};
  \node[p-free] at (0.00,-6.30) {PB2};
  \node[p-dig] at (0.62,-6.30) {PE7};
  \node[d, anchor=west] at (0.96,-6.30) {HX711 DOUT célula 2};
  \node[p-dig] at (0.00,-6.72) {PE8};
  \node[d, anchor=east] at (-0.50,-6.72) {HX711 DOUT célula 3};
  \node[p-dig] at (0.62,-6.72) {PE9};
  \node[d, anchor=west] at (0.96,-6.72) {HX711 DOUT célula 4};
  \node[p-dig] at (0.00,-7.14) {PE10};
  \node[d, anchor=east] at (-0.50,-7.14) {HX711 DOUT célula 5};
  \node[p-dig] at (0.62,-7.14) {PE11};
  \node[d, anchor=west] at (0.96,-7.14) {HX711 DOUT célula 6};
  \node[p-free] at (0.00,-7.56) {PE12};
  \node[p-free] at (0.62,-7.56) {PE13};
  \node[p-free] at (0.00,-7.98) {PE14};
  \node[p-dig] at (0.62,-7.98) {PE15};
  \node[d, anchor=west] at (0.96,-7.98) {LED sistema pronto};
  \node[p-com] at (0.00,-8.40) {PB10};
  \node[d, anchor=east] at (-0.50,-8.40) {I2C2 SCL (EEPROM)};
  \node[p-com] at (0.62,-8.40) {PB11};
  \node[d, anchor=west] at (0.96,-8.40) {I2C2 SDA (EEPROM)};
  \node[p-pwr] at (0.00,-8.82) {GND};
  \node[p-pwr] at (0.62,-8.82) {3V3};
  \node[p-pwr] at (0.00,-9.24) {GND};
  \node[p-pwr] at (0.62,-9.24) {3V3};
  \node[p-pwr] at (0.00,-9.66) {GND};
  \node[p-pwr] at (0.62,-9.66) {3V3};
  \node[p-pwr] at (8.00,0.00) {3V3};
  \node[p-pwr] at (8.62,0.00) {GND};
  \node[p-pwr] at (8.00,-0.42) {VBAT};
  \node[p-pwr] at (8.62,-0.42) {GND};
  \node[p-free] at (8.00,-0.84) {PB9};
  \node[p-free] at (8.62,-0.84) {PB8};
  \node[p-free] at (8.00,-1.26) {PB7};
  \node[p-dig] at (8.62,-1.26) {PB6};
  \node[d, anchor=west] at (8.96,-1.26) {PWM do ESC (TIM4 CH1)};
  \node[p-free] at (8.00,-1.68) {PB5};
  \node[p-dig] at (8.62,-1.68) {PB4};
  \node[d, anchor=west] at (8.96,-1.68) {LED verde};
  \node[p-free] at (8.00,-2.10) {PB3};
  \node[p-free] at (8.62,-2.10) {PD7};
  \node[p-free] at (8.00,-2.52) {PD6};
  \node[p-free] at (8.62,-2.52) {PD5};
  \node[p-free] at (8.00,-2.94) {PD4};
  \node[p-free] at (8.62,-2.94) {PD3};
  \node[p-free] at (8.00,-3.36) {PD2};
  \node[p-seg] at (8.62,-3.36) {PD1};
  \node[d, anchor=west] at (8.96,-3.36) {Botão de emergência};
  \node[p-free] at (8.00,-3.78) {PD0};
  \node[p-free] at (8.62,-3.78) {PC12};
  \node[p-free] at (8.00,-4.20) {PC11};
  \node[p-free] at (8.62,-4.20) {PC10};
  \node[p-dig] at (8.00,-4.62) {PA15};
  \node[d, anchor=east] at (7.66,-4.62) {LED vermelho};
  \node[p-free] at (8.62,-4.62) {PA14};
  \node[p-free] at (8.00,-5.04) {PA13};
  \node[p-free] at (8.62,-5.04) {PA12};
  \node[p-free] at (8.00,-5.46) {PA11};
  \node[p-com] at (8.62,-5.46) {PA10};
  \node[d, anchor=west] at (8.96,-5.46) {USART1 RX};
  \node[p-com] at (8.00,-5.88) {PA9};
  \node[d, anchor=east] at (7.66,-5.88) {USART1 TX};
  \node[p-free] at (8.62,-5.88) {PA8};
  \node[p-com] at (8.00,-6.30) {PC9};
  \node[d, anchor=east] at (7.66,-6.30) {Ethernet TX$-$};
  \node[p-com] at (8.62,-6.30) {PC8};
  \node[d, anchor=west] at (8.96,-6.30) {Ethernet TX+};
  \node[p-com] at (8.00,-6.72) {PC7};
  \node[d, anchor=east] at (7.66,-6.72) {Ethernet RX$-$};
  \node[p-com] at (8.62,-6.72) {PC6};
  \node[d, anchor=west] at (8.96,-6.72) {Ethernet RX+};
  \node[p-free] at (8.00,-7.14) {PD15};
  \node[p-free] at (8.62,-7.14) {PD14};
  \node[p-free] at (8.00,-7.56) {PD13};
  \node[p-free] at (8.62,-7.56) {PD12};
  \node[p-free] at (8.00,-7.98) {PD11};
  \node[p-free] at (8.62,-7.98) {PD10};
  \node[p-free] at (8.00,-8.40) {PD9};
  \node[p-free] at (8.62,-8.40) {PD8};
  \node[p-free] at (8.00,-8.82) {PB15};
  \node[p-free] at (8.62,-8.82) {PB14};
  \node[p-seg] at (8.00,-9.24) {PB13};
  \node[d, anchor=east] at (7.66,-9.24) {Relé de pré-carga};
  \node[p-seg] at (8.62,-9.24) {PB12};
  \node[d, anchor=west] at (8.96,-9.24) {Relé de potência};
  \node[p-pwr] at (8.00,-9.66) {5V};
  \node[p-pwr] at (8.62,-9.66) {GND};
  % legenda (centralizada sob a placa)
  \begin{scope}[local bounding box=leg, shift={(0,-10.75)}]
    \node[p-nm, minimum width=0.4cm] (l0) at (0,0) {};
    \node[d, right=2pt of l0] (t0) {ADC não montado};
    \node[p-dig, minimum width=0.4cm, right=10pt of t0] (l1)  {};
    \node[d, right=2pt of l1] (t1) {E/S digital};
    \node[p-seg, minimum width=0.4cm, right=10pt of t1] (l2)  {};
    \node[d, right=2pt of l2] (t2) {Segurança e relés};
    \node[p-com, minimum width=0.4cm, right=10pt of t2] (l3)  {};
    \node[d, right=2pt of l3] (t3) {Comunicação};
    \node[p-pwr, minimum width=0.4cm, right=10pt of t3] (l4)  {};
    \node[d, right=2pt of l4] (t4) {Alimentação};
    \node[p-res, minimum width=0.4cm, right=10pt of t4] (l5)  {};
    \node[d, right=2pt of l5] (t5) {Reserva};
  \end{scope}
\end{tikzpicture}
}

    Fonte: Autor (2026)
\end{figure}

\subsection{\textit{Firmware}}

A \autoref{fig:blocos_siba} situa as camadas da \autoref{sec:vcore}
dentro do \textit{firmware} da placa. A passagem de \textit{System} para
\textit{Vcore - OS} é sempre mediada pelo OSSI, cujo mecanismo de fila de
trabalho única já foi detalhado na \autoref{fig:arquitetura_fila_unica}. O
\textit{shell} de diagnóstico (\texttt{Test\_Shell.c}) é a única exceção a
esse barramento único: por rodar já dentro do \textit{firmware} e não
depender da rede, ele fala direto com os gerenciadores de domínio, evitando
serializar e desserializar um comando que nunca sai da placa.

\begin{figure}[H]
    \centering
    \caption{Camadas do \textit{firmware} da placa de medição e módulos de cada uma}
    \label{fig:blocos_siba}
    \begin{tikzpicture}[
        font=\small,
        layer/.style={draw, thick, rounded corners=3pt, minimum width=11.6cm},
        mod/.style={draw, rounded corners=2pt, align=center, minimum height=0.8cm, minimum width=1.9cm, fill=white},
        title/.style={font=\bfseries\small},
        arrow/.style={-{Latex}, thick}
    ]
        \node[layer, fill=blue!15, minimum height=1.0cm] (app) at (0,4.3) {\textbf{Application}};

        \node[layer, fill=red!12, minimum height=2.4cm] (sys) at (0,1.85) {};
        \node[title] at (sys.north) [yshift=-0.3cm] {System};
        \node[draw, rounded corners=2pt, fill=yellow!20, minimum width=11.0cm, minimum height=1.55cm] at (0,1.6) {};
        \node[mod] at (-4.4,1.6)  {\textit{Battery}\\\textit{Manager}};
        \node[mod] at (-2.2,1.6)  {\textit{Propeller}\\\textit{Manager}};
        \node[mod] at (0,1.6)     {\textit{Observer}\\\textit{Manager}};
        \node[mod] at (2.2,1.6)   {\textit{Security}\\\textit{Manager}};
        \node[mod] at (4.4,1.6)   {\textit{Log}\\\textit{Manager}};

        \node[layer, fill=magenta!12, minimum height=1.7cm] (os) at (0,-0.85) {};
        \node[title] at (os.north) [yshift=-0.3cm] {Vcore - OS};
        \node[mod] at (-4.4,-1.1) {OSSI};
        \node[mod] at (-2.2,-1.1) {NCI};
        \node[mod] at (0,-1.1)    {HAL};
        \node[mod] at (2.2,-1.1)  {CMD\_Core};
        \node[mod] at (4.4,-1.1)  {NVM};

        \draw[arrow] (app.south) -- (sys.north);
        \draw[arrow] (sys.south) -- (os.north) node[midway, right]{via OSSI};
    \end{tikzpicture}

    Fonte: Autor (2026)
\end{figure}

A camada \textit{System} concentra a lógica de negócios do sistema: os
gerenciadores de domínio listados a seguir, sem nenhuma \textit{thread}
própria: eles são chamados diretamente pela função de despacho da fila de
trabalho única do OSSI (\autoref{sec:vcore}), nunca ao mesmo
tempo. A Figura \ref{fig:Metodologia_SystemManager_Siba} detalha a composição desses gerenciadores:

\begin{itemize}
    \item \textbf{\textit{Propeller Manager}:} Gerencia os parâmetros da hélice: o acelerador solicitado ao ESC, limitado conforme a corrente.
    \item \textbf{\textit{Observer Manager}:} Executa o observador por modos deslizantes com o PLL (\autoref{sec:same_firmware}), estimando rotação e torque a partir das tensões e correntes de fase.
    \item \textbf{\textit{Battery Manager}:} Comanda a energização do conjunto, fechando o relé de pré-carga e depois o de potência, e o desligamento, abrindo os dois.
    \item \textbf{\textit{Security Manager}:} Verifica as tensões e correntes de fase (sobrecorrente e leitura fora da faixa), o estado do observador (divergência e rotação implausível) e a parada de emergência, desligando o motor em situação de risco.
    \item \textbf{\textit{Log Manager}:} Responsável pela telemetria do sistema. Emprega a técnica de agregação de dados (\textit{batching}) para empacotar respostas de comandos e leituras de sensores em um único vetor, otimizando o uso da banda de rede ao enviar os dados de volta para a estação de controle.
\end{itemize}

\begin{figure}[H]
    \centering
    \caption{Diagrama de blocos do \textit{System}}
    \label{fig:Metodologia_SystemManager_Siba}
    \begin{tikzpicture}[font=\small, >={Latex},
        camada/.style={draw, thick, rounded corners=3pt, minimum width=12cm},
        ger/.style={draw, rounded corners=2pt, fill=white, minimum width=3.2cm, minimum height=0.9cm, align=center}]
        \node[camada, fill=blue!15, minimum height=1.1cm] (app) at (0,4.6) {\textit{Application}};
        \node[camada, fill=red!12, minimum height=3.6cm] (sys) at (0,1.5) {};
        \node[anchor=north west, font=\small\bfseries] at (sys.north west) {\textit{System}};
        \node[draw, rounded corners=2pt, fill=yellow!20, minimum width=11cm, minimum height=2.6cm] (sm) at (0,1.3) {};
        \node[anchor=north west, font=\footnotesize] at (sm.north west) {Gerenciadores de domínio};
        \node[ger] at (-3.8,1.55) {\textit{Propeller}\\\textit{Manager}};
        \node[ger] at (0,1.55)    {\textit{Battery}\\\textit{Manager}};
        \node[ger] at (3.8,1.55)  {\textit{Security}\\\textit{Manager}};
        \node[ger] at (-1.9,0.45) {\textit{Observer Manager}};
        \node[ger] at (1.9,0.45)  {\textit{Log Manager}};
        \node[camada, fill=magenta!12, minimum height=1.1cm] (os) at (0,-1.6) {\textit{Vcore - OS}: OSSI (interface entre o sistema operacional e o \textit{System})};
        \draw[<->, very thick] (sys.north) -- node[right, font=\footnotesize]{requisições da aplicação} (app.south);
        \draw[<->, very thick] (os.north) -- (sm.south);
    \end{tikzpicture}

    Fonte: Autor (2026)
\end{figure}

A \autoref{fig:sequencia_siba} resume os dois caminhos que convergem na
fila do OSSI, um valor medido virando telemetria e um comando recebido
virando ação, que nunca são executados ao mesmo tempo.

\begin{figure}[H]
    \centering
    \caption{Sequência de processamento de um sensor e de um comando na placa de medição}
    \label{fig:sequencia_siba}
    \resizebox{\textwidth}{!}{%
    \begin{tikzpicture}[
        font=\scriptsize,
        actor/.style={draw, rounded corners, align=center, minimum width=2.1cm, minimum height=0.7cm},
        lifeline/.style={dashed, gray},
        msg/.style={-{Latex}, thick}
    ]
        \node[actor] (hal) at (0,0)    {HAL};
        \node[actor] (q)   at (2.3,0)  {Fila};
        \node[actor] (wq)  at (4.6,0)  {Fila única\\\scriptsize(OSSI)};
        \node[actor] (cmd) at (6.9,0)  {CMD\_Core};
        \node[actor] (mg)  at (9.2,0)  {Gerenciadores};
        \node[actor] (nci) at (11.5,0) {NCI};
        \node[actor] (pc)  at (13.8,0) {PC};

        \foreach \n in {hal,q,wq,cmd,mg,nci,pc}
            \draw[lifeline] (\n.south) -- ++(0,-12.2);

        \draw[msg] (0,-1.0)  -- (2.3,-1.0)  node[midway, above]{\texttt{Submit sensor}};
        \draw[msg] (2.3,-2.0) -- (4.6,-2.0) node[midway, above]{\textit{dispatch}};
        \draw[msg] (4.6,-3.0) -- (9.2,-3.0) node[midway, above]{Log, Sec, Batt, Prop};
        \node[gray] at (9.2,-4.0) {\textit{temporizador 100 ms}};
        \draw[msg] (9.2,-5.0) -- (11.5,-5.0) node[midway, above]{\texttt{TELEMETRY\_LOG}};
        \draw[msg] (11.5,-6.0) -- (13.8,-6.0) node[midway, above]{UDP 5555};

        \draw[msg] (13.8,-7.4) -- (11.5,-7.4) node[midway, above]{UDP 5000, \texttt{CMD}};
        \draw[msg] (11.5,-8.4) -- (2.3,-8.4)  node[midway, above]{\texttt{Submit CMD}};
        \draw[msg] (2.3,-9.4) -- (4.6,-9.4)   node[midway, above]{\textit{dispatch}};
        \draw[msg] (4.6,-10.4) -- (6.9,-10.4) node[midway, above]{\texttt{CMD\_Core\_Dispatch}};
        \draw[msg] (6.9,-11.3) -- (0,-11.3)   node[midway, above, font=\tiny]{ação direta (ESC, calibração)};
        \draw[msg] (6.9,-12.0) -- (9.2,-12.0) node[midway, above, font=\tiny]{ação via gerenciador};
    \end{tikzpicture}%
    }

    Fonte: Autor (2026)
\end{figure}

Os comandos seguem o formato genérico da \autoref{fig:yaml_example_config};
os disponíveis na placa estão na \autoref{tab:comandos_placa}. No sentido
inverso, a telemetria reúne todas as condições de alarme em um único campo de
32 bits, com um bit por condição.

\begin{table}[H]
    \centering
    \caption{Comandos aceitos pela placa de medição}
    \label{tab:comandos_placa}
    \small
    \begin{tabular}{lp{10.5cm}}
        \toprule
        \textbf{Comando} & \textbf{Função} \\
        \midrule
        Partida     & Arma o conjunto (pré-carga, potência e controlador) ou o desliga \\
        Parada      & Parada de emergência: desarma o controlador do motor \\
        Reinício    & Reinicia a placa \\
        Sensor      & Calibração dos canais e parâmetros do motor e do observador \\
        Acelerador  & Define o acelerador (0 a \SI{100}{\percent}), limitado pela corrente \\
        Alimentação & Aciona os relés de alimentação \\
        \bottomrule
    \end{tabular}

    \vspace{2pt}
    {\small Fonte: Autor (2026)}
\end{table}

Por fim, a camada de \textbf{\textit{Application}} não fica no microcontrolador:
ela é a aplicação de supervisão no computador (\autoref{sec:interface}), que
envia comandos e recebe a telemetria pela rede. Essa separação reduz o
consumo de memória do \textit{firmware}, que mantém apenas o que precisa
responder em tempo real, e permite mudar os ensaios sem regravar a placa.
